\begin{definition}[$\psi_{m,Q,\epsilon,C}$]
  Let $E$ be a metric space, $\gamma \in \Theta(E)$ (\noderef{Curve}), $m \in \mathbb{N}$,
  $Q \subset E$ a finite set, and $\epsilon, C \in \mathbb{R}$. Write $s_i := i \cdot
  \length(\gamma)/m$ for $i = 0,\dotsc,m$. Define
  \[
    \psi_{m,Q,\epsilon,C}(\gamma) := C\sum_{i=1}^m \min \left\{ 2C \frac{\length(\gamma)}{m}, \,
      2\epsilon+ 2C\bigl(d(\gamma(s_i),Q) + d(\gamma(s_{i-1}),Q)\bigr) \right\}
      + 2C^2\frac{\length(\gamma)^2}{m}
    ,
  \]
  where $d(x,Q) = \inf_{q \in Q} d(x,q)$ (Mathlib's \texttt{Metric.infDist}, with the convention
  $d(x,\emptyset) = 0$). This is exactly the first display of paper Lemma~A.2 (paper.tex line
  1718--1721), transcribed with $Q$ ranging over arbitrary finite subsets of $E$ (rather than of a
  fixed countable dense $E_{\mathbb{Q}} \subset E$) and $\epsilon, C$ ranging over all of
  $\mathbb{R}$ (rather than only $\mathbb{Q}_{>0}$): the paper's restriction to rational $\epsilon,
  C$ and to $Q \subset E_{\mathbb{Q}}$ (paper.tex line 1292) is a countability device used only for
  the Strong-Law-of-Large-Numbers step of Section~4.2 (paper.tex line 1291--1295), not part of the
  mathematical content of $\psi_{m,Q,\epsilon,C}$'s definition or of Lemma~A.2 / Corollary~A.6, both
  of which hold verbatim for arbitrary real $\epsilon, C$ and arbitrary finite $Q \subset E$; using
  the wider (paper-justified) quantification here is therefore a strengthening the paper itself
  licenses, not an unexplained one.
\end{definition}
